\section*{\centering{RESUMO}}
\noindent O processamento e a auditoria de documentos da Escrituração Fiscal
Digital de ICMS e IPI (EFD ICMS IPI) enfrentam gargalos operacionais no setor
público e privado. A estrutura em texto plano hierárquico desses arquivos, a
alta volatilidade dos leiautes
legais do SPED e a ausência de uma especificação oficial legível por máquina
tornam os sistemas tradicionais baseados em código fixo inflexíveis e
propensos a falhas. Para solucionar esse problema, o presente
trabalho desenvolve um \textit{pipeline} de ingestão dinâmico e orientado 
metadados.
A arquitetura elaborada extrai as regras de negócio da documentação fiscal e
as
estrutura em arquivos de configuração JSON, desacoplando o processamento
do código-fonte. A solução é composta por um Orquestrador que distribui
a carga de trabalho em lotes e gerencia a execução paralela de trabalhadores
dinâmicos em memória, responsáveis por interpretar a hierarquia do texto e
realizar a persistência estruturada em um banco de dados relacional
\textit{PostgreSQL}. A validação empírica atestou zero perda informacional
e a manutenção da integridade referencial por meio do teste de reconstrução
reversa dos dados originais. A solução também demonstrou alta resiliência
no tratamento de arquivos corrompidos e garantiu a atomicidade das operações
no banco. Conclui-se que a abordagem por metadados confere alta flexibilidade
às mudanças da legislação tributária, e as limitações encontradas oferecem
largo espaço para desenvolvimento e expansões futuras da ferramenta.

\vspace{\baselineskip}

\noindent {\bf Palavras-chave:} EFD ICMS IPI.
SPED.
Engenharia de Dados.
Ingestão de Dados.
Arquitetura Orientada a Metadados.
Banco de Dados Relacional.